\section{Giới thiệu}

\subsection*{Bối cảnh}
Với sự phổ biến của các thiết bị nhúng giá rẻ như ESP32 và nhu cầu ngày càng tăng về giám sát và điều khiển từ xa trong các hệ thống nhà thông minh, việc xây dựng một nền tảng tích hợp để thu thập dữ liệu cảm biến, lưu trữ lịch sử, hiển thị trực quan và điều khiển thiết bị trở nên thiết yếu. Đồ án \textit{Simplehome} nhằm thiết kế và triển khai một hệ thống IoT kết nối các board Yolo Bit (ESP32) với một nền tảng server/dashboard hiện đại, với các yêu cầu về độ tin cậy, bảo mật và khả năng mở rộng.

\subsection*{Mục tiêu chính}
Đồ án đặt ra các mục tiêu chính sau:
\begin{itemize}
    \item Thiết kế và triển khai kiến trúc phần mềm cho hệ thống IoT hỗ trợ thu thập dữ liệu thời gian thực, lưu trữ chuỗi thời gian và truy vấn hiệu quả.
    \item Xây dựng cầu nối giữa giao thức nhắn tin nhẹ MQTT của các thiết bị cảm biến và API REST phục vụ giao diện người dùng (frontend).
    \item Đảm bảo các tính chất quan trọng: xác thực đa người dùng, ghi nhận lịch sử hành động (audit trail), bảo mật dữ liệu và triển khai an toàn bằng containerization.
    \item Đánh giá hiệu năng cơ bản, độ trễ của hệ thống và tính ổn định khi mở rộng số lượng thiết bị.
\end{itemize}

\subsection*{Phạm vi của đồ án}
Đồ án bao gồm:
\begin{itemize}
    \item Một backend Node.js đóng vai trò MQTT client, chuyển tiếp dữ liệu cảm biến vào cơ sở dữ liệu PostgreSQL và cung cấp RESTful API cho các yêu cầu từ frontend.
    \item Một frontend Next.js cung cấp giao diện người dùng: đăng nhập, bảng điều khiển thời gian thực, biểu đồ lịch sử dữ liệu cảm biến và các widget điều khiển thiết bị.
    \item Triển khai toàn bộ hệ thống bằng Docker Compose để đơn giản hóa cài đặt, cấu hình và tái tạo môi trường.
    \item Mô tả chi tiết kiến trúc hệ thống, luồng dữ liệu, giao thức truyền thông và các biện pháp bảo mật.
\end{itemize}

\subsection*{Cấu trúc báo cáo}
Báo cáo được tổ chức thành sáu phần chính:
\begin{enumerate}
    \item \textbf{Giới thiệu} — trình bày bối cảnh, mục tiêu và phạm vi của đồ án.
    \item \textbf{Cơ sở lý thuyết} — giới thiệu các khái niệm, công nghệ và tiêu chuẩn liên quan (IoT, MQTT, Backend/Frontend, Container, v.v.).
    \item \textbf{Phân tích và Thiết kế hệ thống} — mô tả chi tiết kiến trúc hệ thống, lớp ứng dụng, lưu trữ dữ liệu, luồng tương tác và giao thức truyền thông.
    \item \textbf{Quá trình Cài đặt và Triển khai} — trình bày các công nghệ sử dụng, môi trường phát triển, và các bước thực hiện triển khai hệ thống.
    \item \textbf{Kết quả thực nghiệm và Đánh giá} — đánh giá hiệu năng, độ trễ hệ thống, khả năng mở rộng và tính ổn định.
    \item \textbf{Kết luận và Hướng phát triển} — tóm tắt những đóng góp chính, hạn chế hiện tại và các hướng cải tiến trong tương lai.
\end{enumerate}